\chapter{I ritardi}
\label{chap:ritardi}

Oltre alle dissonanze "strutturali" contenute negli accordi di settima, l'armonia usa dissonanze melodiche, cioè note estranee all'accordo che creano una tensione momentanea prima di risolvere. Il \textbf{ritardo} è una delle principali figure di questa categoria.

\section{Definizione e meccanismo del ritardo}
\label{sec:definizione_ritardo}

Il ritardo è una dissonanza che si manifesta sul tempo forte di una misura (o su una parte forte di un tempo), ritardando, appunto, l'arrivo di una nota reale dell'accordo. Si svolge in tre fasi, da rispettare con rigore.

\begin{definizione}[Le tre fasi del Ritardo]
	\label{def:fasi_ritardo}
	Un ritardo si articola sempre in tre momenti distinti:
	\begin{enumerate}
		\item \textbf{Preparazione:} La nota che diventerà dissonante viene introdotta come nota consonante nell'accordo precedente, preferibilmente nella stessa voce e alla stessa altezza.
		\item \textbf{Percussione:} L'accordo cambia, ma la nota preparata viene tenuta (legata), diventando così dissonante rispetto al nuovo contesto armonico. Questa fase deve avvenire su un tempo forte.
		\item \textbf{Risoluzione:} La nota dissonante scende di grado congiunto su una nota reale dell'accordo.
	\end{enumerate}
\end{definizione}

\section{Il ritardo 4-3}
\label{sec:ritardo_4_3}

Il ritardo 4-3 è il più usato. Ritarda l'arrivo della terza dell'accordo: sulla dominante ritarda la sensibile (Si in Do maggiore), sulla tonica ritarda la terza dell'accordo (Mi in Do maggiore).

\subsection{Il ritardo nella cadenza perfetta}
\label{subsec:ritardo_cadenza}

Il ritardo 4-3 si usa soprattutto nella cadenza perfetta (V-I), dove può amplificare la tensione della dominante e ritardare la risoluzione sulla tonica. Come per la settima di dominante, esistono due forme principali di realizzazione.

\begin{esempio}[Forma A: \sigla{V}{7}{} completo → I incompleto con ritardo 4-3 in Do Maggiore]
	\label{ex:ritardo_completo_maggiore}
	Questa progressione utilizza tutti e quattro i suoni sulla dominante (compresa la settima) e risolve su un accordo di tonica incompleto (senza quinta), con raddoppio della fondamentale e ritardo della terza.
	
	\musica{es_cap10_01}
	
	La preparazione del ritardo avviene nell'accordo di IV: il Fa (nel contralto) è nota reale dell'accordo. Con l'arrivo del \sigla{V}{7}{} (Sol al basso) il Fa è tenuto (legato) e diventa la settima; il soprano scende alla sensibile Si e il tenore salta al Re (quinta). Nella seconda battuta, sul basso Do, il Fa è percosso come ritardo 4 e risolve su Mi (3), mentre la sensibile sale a Do e il Re scende a Do.
\end{esempio}

\begin{esempio}[Forma B: \sigla{V}{7}{} incompleto → I completo con ritardo 4-3 in Do Maggiore]
	\label{ex:ritardo_incompleto_maggiore}
	Questa progressione omette la quinta dalla dominante, permettendo di risolvere su un accordo di tonica completo con ritardo della terza.
	
	\musica{es_cap10_02}
	
	In questa versione la dominante è incompleta (senza Re, la quinta) e la fondamentale è raddoppiata; la settima (Fa) è tenuta dal contralto e diventa il ritardo. L'accordo di tonica risultante è completo, con fondamentale, terza (ritardata) e quinta.
\end{esempio}

\subsection{Il ritardo nel modo minore}
\label{subsec:ritardo_minore}

Nel modo minore, il ritardo 4-3 segue gli stessi principi, ma si adatta alle caratteristiche della scala minore armonica, in particolare all'uso della sensibile innalzata nel \sigla{V}{7}{}.

\begin{esempio}[Forma A: \sigla{V}{7}{} completo → I incompleto con ritardo 4-3 in La minore]
	\label{ex:ritardo_completo_minore}
	Progressione con \sigla{V}{7}{} completo che risolve su I incompleto con ritardo.
	
	\musica{es_cap10_03}
\end{esempio}

\begin{esempio}[Forma B: \sigla{V}{7}{} incompleto → I completo con ritardo 4-3 in La minore]
	\label{ex:ritardo_incompleto_minore}
	Progressione con \sigla{V}{7}{} incompleto che risolve su I completo con ritardo.
	
	\musica{es_cap10_04}
\end{esempio}

\begin{regola}[Cifratura del Ritardo]
	\label{reg:cifratura_ritardo}
	Nel basso figurato, il ritardo 4-3 si indica scrivendo in sequenza \texttt{4 3}. Il numero 4 rappresenta la dissonanza (la quarta rispetto al basso), mentre il 3 indica la sua risoluzione (la terza dell'accordo).
\end{regola}

Il ritardo 4-3 arricchisce la progressione cadenzale con uno slancio melodico nelle voci superiori, mentre il basso resta di sostegno. La scelta tra forma completa e incompleta dipende dal contesto e dal volere o no tutte le note dell'accordo di arrivo.

\section{Il ritardo 4-3 con il ritardo nel soprano e nel contralto}
\label{sec:ritardo_4_3_voci}

Il ritardo 4-3 sulla dominante può essere affidato a qualunque voce superiore. Qui il ritardo cade sulla dominante, la nota ritardata (il Do in Do maggiore, il La in La minore) è la fondamentale della tonica che precede, è \textbf{tenuta} (legata) sul basso Sol (o Mi) e scende di grado sulla terza dell'accordo, cioè sulla sensibile. La sensibile non compare in nessun'altra voce durante il ritardo.

\begin{esempio}[Ritardo 4-3 nel soprano]
	\label{ex:ritardo_4_3_s}
	\musica{es_cap10_4_3_s}
\end{esempio}

\begin{esempio}[Ritardo 4-3 nel contralto]
	\label{ex:ritardo_4_3_a}
	\musica{es_cap10_4_3_a}
\end{esempio}

\section{Gli altri ritardi: 9-8, 9-3, 7-6, 7-\texorpdfstring{$\sharp$}{\#}6}
\label{sec:altri_ritardi}

La cifra \textbf{X--Y} indica l'intervallo del ritardo sopra il basso (X) e quello della risoluzione (Y). Oltre al 4-3 si usano quattro ritardi, di cui conviene conoscere la preparazione, il ritardo e l'accordo di risoluzione.

\begin{center}
\begin{tabular}{|l|l|l|l|}
\hline
Ritardo & Preparazione & Nota ritardata & Risoluzione \\
\hline
9-8 & V (quinta del V) & nona sul basso & fondamentale: tonica in stato fondamentale \\
9-3 & IV, iv (terza) & nona sul V & terza di una tonica in primo rivolto (il basso scende di terza) \\
7-6 & V (quinta del V) & settima sul basso & sesta: tonica in primo rivolto \\
7-$\sharp$6 & I (fondamentale) & settima sul II grado & sesta \emph{maggiore}: \sigla{VII}{6}{} (triade diminuita in primo rivolto) \\
\hline
\end{tabular}
\end{center}

\begin{regola}[Condotta del ritardo]
	\label{reg:condotta_ritardo}
	\begin{itemize}
		\item La nota ritardata è \textbf{la stessa} nella preparazione e nel ritardo (legatura di valore).
		\item Il ritardo cade sul tempo forte, la risoluzione sul debole, \textbf{sempre per grado discendente}.
		\item Durante il ritardo la \textbf{nota di risoluzione non compare in nessuna voce superiore}, per evitare la seconda (o nona) aspra. Fa eccezione la risoluzione sulla stessa nota del basso, come nel 9-8 (il Do è già nel basso).
		\item La \textbf{sensibile} del V sale alla tonica. Quando la tonica è proprio la nota di risoluzione del ritardo (9-8, 7-6) non può salire durante il ritardo: se la sensibile è in una \textbf{voce interna} può scendere di terza sulla quinta della tonica (eccezione di Bach, \ref{oss:eccezione_bach}); non può mai scendere nel soprano.
		\item Si evitano le quinte e le ottave parallele.
	\end{itemize}
\end{regola}

\begin{esempio}[Ritardo 9-8 nel soprano]
	\label{ex:ritardo_9_8_s}
	La nona Re (quinta del V) è tenuta sul Do del basso e scende al Do. La sensibile, nel tenore, cade di terza sul Sol (Do maggiore) o sul Mi (La minore).
	\musica{es_cap10_9_8_s}
\end{esempio}

\begin{esempio}[Ritardo 9-8 nel contralto]
	\label{ex:ritardo_9_8_a}
	Il ritardo del contralto in 9-8 è realizzabile solo in Do maggiore e in disposizione stretta. In La minore non esiste una realizzazione corretta a quattro voci: la sensibile non può salire (comparirebbe il La contro il Si ritardato), nel soprano non può cadere, e nel tenore la caduta di terza obbligherebbe a una distanza eccessiva tra le voci.
	\musica{es_cap10_9_8_a}
\end{esempio}

\begin{esempio}[Ritardo 9-3]
	\label{ex:ritardo_9_3}
	Nel passaggio dal V alla tonica in primo rivolto il basso scende di terza (Sol-Mi, Mi-Do). In La minore la nona sopra il Mi è il Fa, una nona minore (cifra $\flat$9).
	Ritardo nel soprano:
	\musica{es_cap10_9_3_s}
	Ritardo nel contralto:
	\musica{es_cap10_9_3_a}
\end{esempio}

\begin{esempio}[Ritardo 7-6]
	\label{ex:ritardo_7_6}
	La settima (Re sul Mi del basso) scende alla sesta (Do): la risoluzione è un accordo di sesta. Nel soprano:
	\musica{es_cap10_7_6_s}
	Il ritardo 7-6 nel contralto non è realizzabile a quattro voci nel rispetto delle spaziature: la sensibile può cadere solo nel tenore e il soprano, per non produrre ottave parallele con il basso, dovrebbe staccarsi dal contralto di più di un'ottava.
\end{esempio}

\begin{esempio}[Ritardo 7-$\sharp$6]
	\label{ex:ritardo_7_6d}
	La settima Do è tenuta sul Re del basso e scende di semitono alla sensibile Si (sesta maggiore): l'accordo di risoluzione è \sigla{VII}{6}{}. La sensibile sale alla tonica, la quinta diminuita Fa scende a Mi. Ritardo nel soprano:
	\musica{es_cap10_7_6d_s}
	Ritardo nel contralto:
	\musica{es_cap10_7_6d_a}
\end{esempio}

\begin{errore}[Raddoppiare o anticipare la nota di risoluzione]
	\label{err:ritardo_risoluzione}
	Scrivere nel ritardo, in un'altra voce, la nota su cui il ritardo risolve (per esempio il Si accanto al Do ritardato del 4-3) produce un urto di seconda minore o nona, annulla l'effetto del ritardo e viene corretto come errore.
\end{errore}

